\subsection{Messverfahren}
Durch mittiges Eintauchen der Plastikkügelchen mit einer Pinzette im PEG-gefüllten Messzylinder wird der Sinkvorgang initiert. Die Fallgeschwindigkeit wird offensichtlicherweise durch Messen einer Fallstrecke und zugehöriger Fallzeit erhoben - NoShitSherlock. Dies wird pro Durchmesser mit 5 Kugeln wiederholt, sodass pro Durchmesser eine Messreihe evaluiert werden kann. \par
Im zweiten Teil, nach Hagen-Poiseuille, wird, in einer äußerst weltbewegenderen Weise als in Versuchsteil 1, nun durch Messen der Durchflussmenge \(V\) (Becherglas mit Skala) pro Zeit \(t\) (Stoppuhr), die zur Auswertung erforderliche Größe des Durchflussstroms \(\dot{V}\) bestimmt. Zusätzlich ist hier, nach \cref{eq:final_eta_hup}, die Bestimmung des an der Kapillare anliegenden Druckes notwendig, was mithilfe einer Höhenmessung des Fluids im Messzylinder erfolgt.
\vspace{2cm}
\begin{figure}[htbp]
	\centering
	\textbf{\textsf{\Large allpapliche Lebenskrise}}\\\vspace{0.3cm}
	
\includegraphics[width=9cm]{sadpablo.jpg}
	\caption[\textit{Meme}: PAP-Lebenskrise\protect\autocite{sadpablo}]{Hier war Platz auf der Seite.\\Fourieroptik war der beste Versuch (rein nach Durchführung und Theorie)\protect\autocite{sadpablo}}
\end{figure}
\newpage
\subsection{Messprotokoll}
\label{sec:Messprotokoll}
\begin{figure}[!h]
    \includegraphics[width=\textwidth]{212.pdf}
    \caption[Messdaten]{fig:Messprotokoll}
\end{figure}
